\section{Weitere Verbesserungen}
\vspace{-0.25cm}
Im Laufe unserer Projektumsetzung stießen wir an einigen Stellen auf Probleme, deren Beseitigung den zeitlichen Rahmen gesprengt hätte.
Dennoch gab es eine Reihe von Ideen für mögliche Lösungen, auf welche wir im folgenden kurz eingehen: 

\vspace{-0.25cm}
\subsubsection*{Synchronisaton}
\vspace{-0.25cm}
Wie in \hyperref[3-2-rendering]{Abschnitt 3.2 Hinzufügen von Objekten zum Renderer} beschrieben, verwenden wir derzeit einen kritischen Bereich,
um das Hinzufügen von Quadern zum Renderer fehlerfrei zu ermöglichen. Über ein Array vordefinierter Größe im Renderer sollte es möglich
sein, diese Operation mit Hilfe der Thread-IDs zu parallelisieren. 

\subsection*{OpenCL}
\vspace{-0.25cm}
\subsubsection*{Speicherverbrauch}
\vspace{-0.25cm}
Ein zentrales Problem der aktuellen OpenCL-Variante ist der Speicherbedarf auf dem Device. Für größere Baumtiefen übersteigt die benötigte Buffergröße den auf der GPU
verfügbaren Speicherplatz, sodass OpenCL mit \texttt{CL\_MEM\_OBJECT\_ALLOCATION\_FAILURE (-4)} abbricht (vgl. \hyperref[openCL-Fehler]{Abschnitt 4.6})\\
Eine Verbesserung könnte durch aufsplitten der Transfers bzw. Speicherobjekte erreicht werden. Statt für eine komplette Ebene einen großen 
Eingabe- und Ausgabebuffer zu reservieren, könnte die Ebene in mehrere Pakete zerlegt werden, die nacheinander verarbeitet werden. 
Dadurch sinkt die maximal benötigte Buffergröße und damit die Wahrscheinlichkeit, dass Allokationen aufgrund fragmentierten oder zu
knappen Device-Speichers fehlschlagen.

\vspace{-0.25cm}
\subsubsection*{Optimierung der Blockgröße und -anzahl}
\vspace{-0.25cm}
Die Optimierung von Blockgröße und Blockanzahl erwies sich als schwierig. Der Hauptgrund ist, dass sich die Anzahl der zu verarbeitenden
Quader pro Ebene verdoppelt, wodurch sich die optimale Verteilung stetig verändert.\\
Als Ansatz zur Verbesserung bietet sich hier eine Pipelinestruktur an. Die Grundidee ist, fortlaufend Pakete von Quadern zu verarbeiten:
Die CPU schickt ein Paket von Quadern an die GPU.
Die GPU erzeugt daraus die neuen Quader (Folgeebene) und schreibt sie in einen Ausgabebuffer.
Während die GPU bereits am nächsten Paket arbeitet, holt die CPU die Ergebnisse des vorherigen Pakets ab --- parallel dazu kann die CPU den Renderer 
befüllen und die bereits generierten Quader zur Darstellung weiterreichen.\\
So entsteht eine Überlappung von Rechenzeit und Transfer-/Verarbeitungszeit, weil die CPU den Renderer befüllt, während die GPU arbeitet.
Im Idealfall reduziert dies Wartezeiten und steigert den Durchsatz des Systems. \\
Ein praktischer Nebeneffekt wäre, dass die Blockgröße durch die Paketgröße automatisch beschränkt würde, was nicht nur die Wahl passender Work-Group-Größen
erleichtert, sondern auch das zuvor beschrieben Problem mit den Speicherlimits reduzieren könnte.

\newpage
\subsection*{Renderer}
Obwohl wir vor allem mit den parallelen Ansätzen theoretisch recht große Bäume generieren können, ist es uns praktisch derzeit nicht möglich, diese
visuell aufzubereiten. Außerdem wäre in Anlehnung an den Fraktal-Charakter der Pythagorasbäume eine Zoom-Option im Renderer von Vorteil.\\
Um die Effizienz unseres Renderers zu verbessern, böte sich Multithreading oder die Verwendung von Hardwarebeschleunigung an. 
